\chapter{Principios básicos de protección contra radiación externa}

Son tres los factores más importantes que deben tomarse en cuenta a fin de que la dosis que se reciba sea mínima: tiempo, distancia y blindaje.

Mantener el menor tiempo a la mayor distancia posible e interponer el mayor blindaje posible entre la fuente de radiación y la persona garantiza que la dosis será mínima.

La exposición sigue la ley del cuadrado inverso; esto es, varía en razón inversa al cuadrado de la distancia. Matemáticamente se expresa como:
\[
X_1r_1^2=X_2r_2^2,
\]
donde $X_1$ es la rapidez de exposición a la distancia $r_1$ y $X_2$ es la rapidez de exposición a la distancia $r_2$.

Así, si se desea calcular la rapidez de exposición a 4 metros de una cierta fuente puntual que está dando $10\ \mathrm{R/h}$ a un metro de distancia:
\begin{align*}
X_2 &= \left(\frac{r_1}{r_2}\right)^2X_1\\
     &= \left(\frac{1}{4}\right)^2(10)\\
     &= 0.0625(10)=0.625\ \mathrm{R/h}\\
     &= 625\ \mathrm{mR/h}.
\end{align*}

Como se observa, alejarse de 1 a 4 metros reduce la rapidez de exposición de $10\ \mathrm{R/h}$ a solo $625\ \mathrm{mR/h}$.

Por otra parte, el blindaje es un método muy eficiente para reducir la rapidez de exposición. Los rayos gamma se atenúan al pasar por un material y el grado de atenuación depende del material en cuestión. Por ejemplo, se requieren 6.2 pulgadas de concreto para reducir la intensidad de radiación a un décimo de su valor original cuando la fuente es Ir-192, mientras que solo se requieren 0.64 pulgadas de plomo para reducirla en la misma proporción.

De esto se deduce que el plomo es mejor blindaje que el concreto.

Al espesor necesario de un cierto material que reduce la intensidad de radiación a la mitad de su valor original se le llama «capa hemirreductora» y es distinta para cada radioisótopo y para cada material.

También se usa el concepto «capa decirreductora» para indicar el espesor necesario para reducir la intensidad a un décimo de su valor original.

El siguiente cuadro muestra los valores de las capas hemirreductoras y decirreductoras para el concreto y el plomo, y para algunos radioisótopos de interés.

\begin{table}[ht]
\centering
\caption{Capas hemirreductoras (C.H.) y decirreductoras (C.D.).}
\label{tab:capas-reductoras}
\begin{tabular}{|l|c|c|c|c|c|c|}
\hline
\textbf{Fuente} & \multicolumn{2}{c|}{\textbf{Iridio-192}} & \multicolumn{2}{c|}{\textbf{Cobalto-60}} & \multicolumn{2}{c|}{\textbf{Cesio-136}} \\
\cline{2-7}
\textbf{Material} & C.H. & C.D. & C.H. & C.D. & C.H. & C.D. \\
\hline
Concreto & 1.9\,in & 6.2\,in & 2.6\,in & 8.6\,in & 2.1\,in & 7.1\,in \\
Plomo    & 0.19\,in & 0.64\,in & 0.49\,in & 1.6\,in & 0.25\,in & 0.84\,in \\
\hline
\end{tabular}
\end{table}

Por ejemplo, si quisiéramos saber el espesor necesario de plomo para blindar una fuente de Ir-192 y reducir su rapidez de exposición de $10\ \mathrm{R/h}$ a $2\ \mathrm{mR/h}$, procederíamos aplicando la siguiente fórmula:
\[
X=\frac{X_0}{2^n},
\]
donde $X$ es la rapidez de exposición, $X_0$ la rapidez de exposición original y $n$ el número de capas hemirreductoras que se necesitan.

Como en este caso interesa conocer $n$, se despeja obteniendo:
\[
n=\frac{\log(X_0/X)}{\log 2}.
\]
Sustituyendo los valores, con $10\ \mathrm{R/h}=10{,}000\ \mathrm{mR/h}$:
\begin{align*}
n &= \frac{\log(10{,}000/2)}{\log 2}\\
  &= \frac{\log(5{,}000)}{\log 2}\\
  &\approx \frac{3.69897}{0.30103}=12.29.
\end{align*}

Es decir, se necesitan aproximadamente 12.29 capas hemirreductoras de plomo. Como cada una de ellas tiene 0.19 pulgadas, el espesor requerido será:
\[
\text{Espesor de plomo necesario}=12.29\times0.19=2.34\ \text{pulgadas}.
\]

\section{Cálculo del tiempo de exposición}

Para calcular el tiempo que puede permanecer una persona en un campo de radiación conocido se utiliza:
\[
t=\frac{H}{X},
\]
donde $H$ es la dosis límite, $X$ la radiación en el área y $t$ el tiempo.

Si se quiere calcular cuánto tiempo puede permanecer un trabajador en un área con $30\ \mathrm{mR/h}$, sin recibir más de $90\ \mathrm{mrem}$:
\[
t=\frac{H}{X}=\frac{90\ \mathrm{mrem}}{30\ \mathrm{mR/h}}=3\ \text{h}.
\]

\section{Cálculos a diferentes distancias}

Para calcular la rapidez de exposición a otra distancia se utiliza la ley del cuadrado inverso:
\[
X_2=\left(\frac{d_1}{d_2}\right)^2X_1.
\]

Por ejemplo, para conocer la rapidez de exposición a 2 metros de una fuente que da $100\ \mathrm{mR/h}$ a 1 metro:
\[
X_2=\left(\frac{1}{2}\right)^2(100)=25\ \mathrm{mR/h}.
\]

Si se quiere calcular cuánto tiempo se puede trabajar a 5 metros de una fuente que da una rapidez de exposición de $2\ \mathrm{R/h}$ a un metro, sin recibir más de $100\ \mathrm{mrem}$, primero se calcula la rapidez de exposición:
\begin{align*}
X_1 &= 2\ \mathrm{R/h}, & d_1&=1\ \mathrm{m},\\
X_2 &= \left(\frac{1}{5}\right)^2(2)=0.08\ \mathrm{R/h}=80\ \mathrm{mR/h}.
\end{align*}

El tiempo para recibir 100 mrem es:
\[
80\ \mathrm{mrem}\longrightarrow 1\ \mathrm{h},\qquad
100\ \mathrm{mrem}\longrightarrow \frac{100}{80}=1.25\ \mathrm{h}=1\ \mathrm{h}\ 15\ \mathrm{min}.
\]

Si en lugar de trabajar 1 hora 15 minutos se trabaja solo 40 minutos, la dosis recibida es:
\[
X=\frac{80\times40}{60}=53.3\ \mathrm{mrem}.
\]

\section{Ejemplos con Ir-192}

La constante gamma de desintegración o emisividad del Ir-192 es:
\[
0.55\ \mathrm{R/h\ a\ 1\ metro\ por\ cada\ curie}.
\]

\begin{center}
\begin{tabular}{cc}
\textbf{Fuente} & \textbf{Rapidez a un metro} \\
\hline
$2\ \mathrm{Ci}$   & $0.55\times2=1.1\ \mathrm{R/h}$ \\
$10\ \mathrm{Ci}$  & $0.55\times10=5.5\ \mathrm{R/h}$ \\
$35\ \mathrm{Ci}$  & $0.55\times35=19.25\ \mathrm{R/h}$ \\
$100\ \mathrm{Ci}$ & $0.55\times100=55\ \mathrm{R/h}$ \\
\end{tabular}
\end{center}

Se plantea calcular cuánto espesor de plomo se necesita para evitar que una fuente de Ir-192 de 53.3 Ci dé más de 100 mR/h a un metro. El original utiliza 53.5 Ci en el cálculo siguiente; se conserva esa discrepancia para revisión de la fuente:
\[
53.5\times0.55=29.425\ \mathrm{R/h}=29{,}425\ \mathrm{mR/h}.
\]

Si 0.2 pulgadas de plomo reducen a la mitad el nivel de radiación, se obtiene:
\begin{table}[ht]
\centering
\caption{Reducción de la radiación con diferentes espesores de plomo.}
\label{tab:plomo-reduccion}
\begin{tabular}{|c|c|}
\hline
\textbf{Espesor} & \textbf{Nivel de radiación} \\
\hline
0      & 29{,}425 mR/h \\
0.2\,in & 14{,}712.5 mR/h \\
0.4\,in & 7{,}356.25 mR/h \\
0.6\,in & 3{,}678.13 mR/h \\
0.8\,in & 1{,}839.1 mR/h \\
1.0\,in & 919.5 mR/h \\
1.2\,in & 459.77 mR/h \\
1.4\,in & 229.88 mR/h \\
1.6\,in & 114.94 mR/h \\
1.8\,in & 57.47 mR/h \\
\hline
\end{tabular}
\end{table}

Para conocer el nivel de radiación a 3 metros de una fuente de Ir-192 de 85 Ci, a un metro se tiene:
\[
0.55\times85=46.75\ \mathrm{R/h}.
\]
Aplicando la ley del cuadrado inverso:
\[
X_2=\left(\frac{1}{3}\right)^2(46.75)=5.19\ \mathrm{R/h}.
\]

\section{Exposición accidental a corta distancia}

A un trabajador se le cae una fuente de 60 Ci y decide tomarla del contenedor con la mano y meterla al contenedor. Si el contenedor está a 10 cm de la fuente, el nivel de radiación en ese punto se calcula a partir de:
\[
0.55\times60=33\ \mathrm{R/h}\quad\text{a un metro}.
\]

A 10 cm, es decir, a 0.1 m:
\[
X_2=\left(\frac{1}{0.1}\right)^2(33)=3{,}300\ \mathrm{R/h}.
\]

Si se tarda 10 segundos en introducir la fuente, la dosis aproximada recibida en la mano sería:
\[
\text{Dosis}=\frac{3{,}300\ \mathrm{R}}{3{,}600\ \mathrm{s}}\times10\ \mathrm{s}=9.17\ \mathrm{R}.
\]

\section{Fuente de Ir-192 con blindaje}

Una fuente de Ir-192 que fue comprada hace 148 días con 120 Ci, y está cubierta con 8 pulgadas de plomo, plantea la pregunta de qué nivel de radiación dará a 2 metros.

Primero se calcula la actividad de la fuente. Ciento cuarenta y ocho días corresponden a dos vidas medias:
\begin{center}
\begin{tabular}{cc}
120 Ci & 0 días \\
60 Ci  & 74 días \\
30 Ci  & 148 días \\
\end{tabular}
\end{center}

Ahora se calcula el nivel de exposición a un metro, sin blindaje:
\[
30\times0.55=16.50\ \mathrm{R/h}.
\]

Con blindaje de plomo:
\begin{center}
\begin{tabular}{cc}
\textbf{Espesor} & \textbf{Nivel de exposición} \\
\hline
0       & 16.5 R/h \\
0.2\,in & 8.25 R/h \\
0.4\,in & 4.12 R/h \\
0.6\,in & 2.06 R/h \\
0.8\,in & 1.03 R/h \\
\end{tabular}
\end{center}

Finalmente, para calcular el nivel de exposición a 2 metros:
\begin{align*}
X_2&=\left(\frac{1}{2}\right)^2(1.03)=0.257\ \mathrm{R/h},\\
X_2&=257\ \mathrm{mR/h}.
\end{align*}
